% Chapter 2
\section{Least Squares: Deriving the Estimator from Scratch}\label{sec:ls}

\subsection{The problem as an optimization problem}

Forget statistics for a moment. We have $n$ points and we want a line (hyperplane) that passes ``as close as possible'' to all of them. ``Close'' must be made precise, and the most natural first guess — the one Legendre and Gauss made — is to measure total misfit by the \emph{sum of squared residuals}:
\begin{equation}\label{eq:rss}
  S(\bbeta) = \sum_{i=1}^{n} \big( y_i - \beta_0 - \beta_1 x_{i1} - \dots - \beta_p x_{ip} \big)^2
            = \norm{\y - \X \bbeta}^2.
\end{equation}
Each term $e_i = y_i - (\X \bbeta)_i$ is the \emph{residual}: the vertical gap between the observation and the fitted hyperplane. Squaring does three useful things at once: it makes every deviation positive, it penalizes large deviations disproportionately, and — crucially, as we will see — it makes the problem \emph{differentiable and convex}, hence exactly solvable. Chapter~\ref{sec:why-squared} interrogates whether this is the ``right'' loss; for now we simply adopt it.

The \textbf{ordinary least squares (OLS) estimator} is
\[
  \bhat \;=\; \arg\min_{\bbeta \in \R^{p+1}} \; \norm{\y - \X \bbeta}^2.
\]

\subsection{The one-variable warm-up: $p = 0$ and $p = 1$}

\textbf{No features ($p=0$): only an intercept.} Minimize $S(\beta_0) = \sum_i (y_i - \beta_0)^2$. Setting the derivative to zero:
\[
  \frac{dS}{d\beta_0} = -2 \sum_{i=1}^n (y_i - \beta_0) = 0
  \quad\Longrightarrow\quad
  \hat{\beta}_0 = \frac{1}{n} \sum_{i=1}^n y_i = \bar{y}.
\]
The least-squares fit with no features is just the mean. Already a theorem is visible: \emph{least squares generalizes the average}. The sample mean is the value that minimizes squared deviation — a fact we will reuse constantly.

\textbf{One feature ($p=1$): the regression line through a scatter plot.} Minimize
\[
  S(\beta_0, \beta_1) = \sum_{i=1}^n (y_i - \beta_0 - \beta_1 x_i)^2.
\]
Take partial derivatives and set them to zero (the ``normal equations'' in disguise):
\[
  \frac{\partial S}{\partial \beta_0} = -2 \sum_i (y_i - \beta_0 - \beta_1 x_i) = 0,
  \qquad
  \frac{\partial S}{\partial \beta_1} = -2 \sum_i x_i (y_i - \beta_0 - \beta_1 x_i) = 0.
\]
From the first equation, $\hat{\beta}_0 = \bar{y} - \hat{\beta}_1 \bar{x}$: the line passes through the centroid $(\bar{x}, \bar{y})$ of the data. Substituting into the second equation and solving:
\begin{equation}\label{eq:slope}
  \hat{\beta}_1 = \frac{\sum_i (x_i - \bar{x})(y_i - \bar{y})}{\sum_i (x_i - \bar{x})^2} = \frac{\Cov_n(x, y)}{\Var_n(x)},
\end{equation}
where $\Cov_n, \Var_n$ denote the sample (dividing by $n$ or $n-1$ — the ratio is the same) covariance and variance.

Read \eqref{eq:slope} slowly: \emph{the slope is the ratio of how $x$ and $y$ vary together to how much $x$ varies alone}. This is the single most quoted formula in applied statistics, and it has fallen out of nothing more than setting a derivative to zero. Nothing in it assumed the errors were Gaussian, independent, or even random — the least-squares machinery is purely geometric optimization.

\begin{example}[Computation by hand]\label{ex:hand}
Take $n = 4$ points $(x, y) = (1, 2), (2, 3), (3, 5), (4, 4)$. Then $\bar{x} = 2.5$, $\bar{y} = 3.5$,
$\sum_i (x_i - \bar{x})(y_i - \bar{y}) = (-1.5)(-1.5) + (-0.5)(-0.5) + (0.5)(1.5) + (1.5)(0.5) = 4$, and $\sum_i (x_i - \bar{x})^2 = 2.25 + 0.25 + 0.25 + 2.25 = 5$, so $\hat{\beta}_1 = 4/5 = 0.8$ and $\hat{\beta}_0 = 3.5 - 0.8 \times 2.5 = 1.5$. The least-squares line is $y = 1.5 + 0.8x$. You can verify the normal equations directly: residuals $e_i = y_i - \hat{y}_i = (2 - 2.3,\; 3 - 3.1,\; 5 - 3.9,\; 4 - 4.7) = (-0.3, -0.1, 1.1, -0.7)$ sum to $0$, and $\sum_i x_i e_i = -0.3 - 0.2 + 3.3 - 2.8 = 0$ — residual orthogonal to the design, as promised.
\end{example}

\subsection{The general case: the normal equations}

Now let $\X \in \R^{n \times (p+1)}$ with $n > p+1$. Expand the objective:
\[
  S(\bbeta) = (\y - \X\bbeta)^\top (\y - \X\bbeta)
            = \y^\top \y - 2 \bbeta^\top \X^\top \y + \bbeta^\top \X^\top \X \bbeta,
\]
using $\y^\top \X \bbeta = \bbeta^\top \X^\top \y$ (both are scalars). Differentiate with respect to the vector $\bbeta$ — using the two elementary identities $\nabla_{\bbeta}(\bm{c}^\top \bbeta) = \bm{c}$ and $\nabla_{\bbeta}(\bbeta^\top \bm{A} \bbeta) = 2 \bm{A} \bbeta$ for symmetric $\bm{A}$:
\begin{equation}\label{eq:grad}
  \nabla_{\bbeta} S = -2 \X^\top \y + 2 \X^\top \X \bbeta.
\end{equation}
Setting the gradient to zero yields the \textbf{normal equations}:
\begin{equation}\label{eq:normal}
  \boxed{\;\X^\top \X \, \bhat \;=\; \X^\top \y\;}
\end{equation}

Provided $\X$ has full column rank (i.e., $\rank(\X) = p+1$, no feature is a linear combination of the others, and $n \geq p+1$), the Gram matrix $\X^\top \X$ is symmetric positive definite hence invertible, and we obtain the celebrated closed form
\begin{equation}\label{eq:ols}
  \boxed{\;\bhat = (\X^\top \X)^{-1} \X^\top \y.\;}
\end{equation}

\begin{remark}[Why ``normal''?]
The equations \eqref{eq:normal} say $\X^\top (\y - \X \bhat) = \zero$: the residual vector is \emph{orthogonal} (normal, in the geometric sense) to every column of $\X$. This is not a coincidence or a curiosity — it is the entire geometric content of least squares, and Chapter~\ref{sec:geometry} is devoted to it.
\end{remark}

\begin{remark}[Why is it a minimum?]
The Hessian $\nabla^2 S = 2 \X^\top \X$ is positive definite when $\X$ has full column rank, so $S$ is strictly convex and the stationary point is the \emph{unique global minimum}. There are no spurious local minima, no initialization issues, no learning rates. Compare that with training a neural network. This exactness is a large part of ``why linear regression.''
\end{remark}

\subsection{Rank deficiency and the pseudoinverse}

What if $\X^\top \X$ is singular — some feature is redundant, or $p \geq n$? The minimizer is no longer unique (a whole flat subspace of $\bbeta$'s attains the minimum), but the \emph{fitted values} $\X \bhat$ still are. The canonical choice is the \textbf{Moore--Penrose pseudoinverse} solution
\[
  \bhat = \X^{+} \y,
\]
which picks, among all minimizers, the one with smallest Euclidean norm. Every serious implementation (R's \texttt{lm}, NumPy's \texttt{lstsq}, scikit-learn's \texttt{LinearRegression}) computes exactly this via the SVD or QR decomposition rather than forming $(\X^\top\X)^{-1}$, for numerical stability. This foreshadows Chapter~\ref{sec:bias-variance}: with redundant or numerous features, constraining the solution norm is not merely a numerical trick, it is a statistical idea (ridge regression).

\subsection{Fitted values, residuals, and $R^2$}

Given $\bhat$, define the \textbf{fitted values} $\hat{\y} = \X \bhat$ and \textbf{residuals} $\bm{e} = \y - \hat{\y}$. Two facts follow immediately from the normal equations:

\begin{enumerate}[label=(\alph*)]
  \item $\X^\top \bm{e} = \zero$: residuals are orthogonal to every column of $\X$. In particular, with an intercept column, $\sum_i e_i = 0$: residuals average to zero.
  \item $\hat{\y}^\top \bm{e} = (\X \bhat)^\top \bm{e} = \bhat^\top \X^\top \bm{e} = 0$: fitted values and residuals are orthogonal.
\end{enumerate}

Fact (b) yields the \textbf{Pythagorean decomposition} (with an intercept):
\[
  \norm{\y - \bar{y} \mathbf{1}}^2
  = \norm{\hat{\y} - \bar{y} \mathbf{1}}^2 + \norm{\bm{e}}^2,
\]
i.e., \emph{total variation = explained variation + unexplained variation}. Dividing through:
\begin{equation}\label{eq:r2}
  R^2 = 1 - \frac{\sum_i e_i^2}{\sum_i (y_i - \bar{y})^2}
      = \frac{\norm{\hat{\y} - \bar{y}\mathbf{1}}^2}{\norm{\y - \bar{y}\mathbf{1}}^2} \in [0, 1].
\end{equation}
$R^2$ is the fraction of variance ``explained'' by the regression, and it is literally the squared cosine of the angle between $\y - \bar{y}\mathbf{1}$ and its projection onto the column space — geometry again.

\begin{keybox}[What we have established so far]
Without a single probabilistic assumption:
\begin{itemize}
  \item the least-squares estimator has the closed form $\bhat = (\X^\top \X)^{-1} \X^\top \y$ (unique iff $\X$ has full column rank);
  \item it is the unique global minimizer (convexity);
  \item residuals are orthogonal to the design matrix (the normal equations);
  \item $R^2$ decomposes variance via the Pythagorean theorem.
\end{itemize}
Everything so far is \emph{deterministic geometry}. Statistics enters when we ask how good $\bhat$ is \emph{as an estimate of the truth} — that is the business of Chapters \ref{sec:gauss-markov} and \ref{sec:probability}.
\end{keybox}
